\chapter{Floating Point Arithmetic}

\section{Why Computers Can't Count Perfectly}

Computers are known to be very precise. Modern processors can perform billions of calculations per second, so it's not intuitive to realize they have this persistent limitation: they cannot store most numbers exactly.

Let us think about what it would actually take to store a number like $\frac{1}{3}$. You
probably know its decimal expansion:
\[
    \frac{1}{3} = 0.3333333\ldots\]
The digits go on forever. Now imagine you are a computer with a fixed amount of memory.
You cannot store infinitely many digits. So you store $0.333$ or $0.3333333$, depending on how much space you have, and you accept that the result is not exactly right. That gap between the number you wanted and the number you
actually stored is called \newterm{approximation error}\index{approximation error}, and
it never fully goes away.

This is not just a problem for fractions with repeating decimals. Irrational numbers like
$\sqrt{2} = 1.41421356\ldots$ and transcendental numbers like $\pi = 3.14159265\ldots$
also have infinite, non-repeating decimal expansions. Every one of them must be
approximated when stored in a computer.

To handle this systematically, mathematicians and engineers developed the
\newterm{floating-point system}\index{floating-point system}. Before we look at how it works, it helps to understand the two core reasons why any finite system will always fall short of
the real number line.

\subsection{The Number Line is Continuous and Unbounded}

Between any two real numbers, no matter how close together, there are infinitely many
others. Pick any two decimals you like. For example, if you picked $1.00$ and $1.01$, you can always find a number between them: $1.005$, $1.0051$, $1.00500001$, and so on without end. This is referred to as \newterm{continuity}\index{continuity}. A computer's number system cannot be continuous. It stores numbers using a fixed string
of digits, which means it can only represent finitely many values. Real numbers have to be rounded to the nearest representable value. 

The real numbers stretch infinitely in both directions. There is no largest real number
and no smallest positive real number. A computer's memory is finite, so its number system must have both a largest and a smallest representable value. Numbers too large to store cause a problem called \newterm{overflow}\index{overflow}. Numbers too small to distinguish from zero cause \newterm{underflow}\index{underflow}. 


% Add example
%test edit

\begin{Exercise}[title={Overflow and Underflow}, label={ex:overflow}]
Suppose a floating-point system can only represent numbers between $10^{-4}$ and $10^{8}$.
\begin{enumerate}
    \item Planck's constant is approximately $6.63 \times 10^{-34}$ J$\cdot$s. Does it
    cause underflow, overflow, or neither in this system?
    \item The speed of light is approximately $3.00 \times 10^{8}$ m/s. Does it
    cause underflow, overflow, or neither in this system?
    \item Give one example of a number that would cause overflow.
\end{enumerate}
\end{Exercise}
\begin{Answer}[ref={ex:overflow}]
\begin{enumerate}
    \item Underflow. $6.63 \times 10^{-34}$ is far below $10^{-4}$.
    \item Neither. $3.00 \times 10^{8}$ is exactly at the upper boundary.
    \item Any number greater than $10^{8}$.
\end{enumerate}
\end{Answer}


So what can we do?

We cannot fix these limitations, because they are baked into the nature of finite computation.
What we can do is design the floating-point system carefully so that the errors it introduces are small, predictable, and manageable. That is exactly what engineers and mathematicians have spent decades working out, and it is what the rest of this chapter is about!

The key insight is that a floating-point number is not just a string of digits. A floating-point number has three parts: a \newterm{sign}\index{sign}, a \newterm{significand}\index{significand}
(sometimes called the mantissa), and an \newterm{exponent}\index{exponent}.


\section{Binary and the Floating-Point System}

In the previous section you saw that computers cannot store most numbers
exactly.
Now the natural question is: how do they store numbers at all?
The answer starts at the hardware level, with the smallest possible piece of
information a computer can hold.

\subsection{Binary System}

Every process inside a CPU comes down to the flow of electrical current through
tiny switches called \newterm{transistors}\index{transistors}.
A transistor is either on (carrying current) or off (not carrying current).
That gives us exactly two states, which we write as $0$ and $1$.
This is why computers use \newterm{binary}\index{binary}: it is the only number
system that hardware can physically implement with a simple on/off switch.

A single binary digit is called a \newterm{bit}\index{bit}.
Eight bits grouped together form a \newterm{byte}\index{byte}.
You have probably seen units like megabyte (MB) or gigabyte (GB) before.
One megabyte is $10^6$ bytes, and one gigabyte is $10^9$ bytes.
A modern smartphone holds tens of billions of bits.

\begin{Exercise}[title={Bits and bytes}, label={ex:bits-bytes}]
  \begin{enumerate}
    \item A file is $4$ megabytes.
          How many bytes is that?
          How many bits?
    \item A hard drive stores $500$ gigabytes.
          How many bits can it hold?
          Write your answer using scientific notation.
    \item A single bit can be $0$ or $1$.
          Two bits together can be $00$, $01$, $10$, or $11$, which are four possible
          values.
          How many distinct values can $8$ bits (one byte) represent?
          How about $64$ bits?
  \end{enumerate}
\end{Exercise}
\begin{Answer}[ref={ex:bits-bytes}]
  \begin{enumerate}
    \item $4 \times 10^6$ bytes;\; $4 \times 10^6 \times 8 = 3.2 \times 10^7$
          bits.
    \item $500 \times 10^9 \times 8 = 4 \times 10^{12}$ bits.
    \item $2^8 = 256$ values for one byte;\;
          $2^{64} \approx 1.8 \times 10^{19}$ values for 64 bits.
  \end{enumerate}
\end{Answer}


\subsection{Floating-Point Number}

Recall that a floating-point number has three parts:
a \newterm{sign}\index{sign}, a
\newterm{significand}\index{significand} (sometimes called the mantissa), and
an \newterm{exponent}\index{exponent}.

Before we write any formulas, let us see why all three parts are necessary
with a base-10 analogy.

Consider the number $-0.001234$.
Written in scientific notation it is $-1.234 \times 10^{-3}$.
Three pieces of information do all the work:
\begin{itemize}
  \item The \textbf{sign} ($-$) tells us the number is negative.
  \item The \textbf{significand} ($1.234$) holds the meaningful digits.
  \item The \textbf{exponent} ($-3$) tells us how large or small the number
        is by shifting the decimal point.
\end{itemize}

Now watch what the exponent does when we fix the significand at $0.1234$ and
vary it in base 10:

\begin{center}
\begin{tabular}{ll}
  \hline
  Float                    & Value \\
  \hline
  $0.1234 \times 10^{-2}$  & $0.001234$ \\
  $0.1234 \times 10^{1}$   & $1.234$ \\
  $0.1234 \times 10^{4}$   & $1234$ \\
  $0.1234 \times 10^{8}$   & $12{,}340{,}000$ \\
  \hline
\end{tabular}
\end{center}

The decimal point \emph{floats}, which means it can sit anywhere, controlled entirely
by the exponent.
That is why these are called floating-point numbers.
The significand always carries the same digits; only the scale changes.

\begin{Exercise}[title={The floating point in action}, label={ex:fp-float}]
  The significand is fixed at $0.5050$ and the base is $10$.
  Fill in the value represented by each float.

  \begin{center}
  \begin{tabular}{ll}
    \hline
    Float                    & Value \\
    \hline
    $0.5050 \times 10^{0}$   & \\[6pt]
    $0.5050 \times 10^{2}$   & \\[6pt]
    $0.5050 \times 10^{-1}$  & \\[6pt]
    $0.5050 \times 10^{5}$   & \\[6pt]
    \hline
  \end{tabular}
  \end{center}

  Now answer: if you want to represent $505{,}000{,}000$ using the same
  significand, what exponent do you need?
\end{Exercise}
\begin{Answer}[ref={ex:fp-float}]
  \begin{itemize}
    \item $0.5050 \times 10^{0} = 0.5050$
    \item $0.5050 \times 10^{2} = 50.50$
    \item $0.5050 \times 10^{-1} = 0.05050$
    \item $0.5050 \times 10^{5} = 50{,}500$
  \end{itemize}
  For $505{,}000{,}000 = 0.5050 \times 10^9$, the exponent needed is $9$.
\end{Answer}

\subsection{Floating-Point System}

A floating-point system is defined by four numbers, written $(\beta, p, m, M)$.
Every float in that system looks like this:
\[
  \pm \;(0.d_1\, d_2\, \cdots\, d_p)_\beta \;\times\; \beta^e
\]

This looks intimidating at first, but each symbol has a simple job.
Let us unpack them one at a time.

\medskip
$\boldsymbol{\beta}$  represents the \newterm{base}\index{base}.
This is the number that the exponent is a power of.
In everyday life we use $\beta = 10$.
Computers use $\beta = 2$ (binary).

\medskip
$\boldsymbol{p}$  represents the \newterm{precision}\index{precision}.
This is how many digits the significand gets to use.
The significand $0.d_1 d_2 \cdots d_p$ has exactly $p$ digits, each chosen
from $\{0, 1, \ldots, \beta - 1\}$.
More digits means more precision, that is, you can represent numbers more
accurately.
Fewer digits means faster arithmetic but more rounding error.

\medskip
$\boldsymbol{e}$  represents the \newterm{exponent}\index{exponent}.
This shifts the significand left or right.
The exponent is not free to be any integer; it is capped between a minimum
$m$ and a maximum $M$.
Those bounds set the range of the system, i.e., how large and how small the
representable numbers can be.

\medskip
$\boldsymbol{m}$ and $\boldsymbol{M}$  represents the \newterm{exponent bounds}\index{exponent bounds}.
The exponent must satisfy $m \leq e \leq M$.
Together they control overflow and underflow.
If a result needs an exponent below $m$, underflow occurs.
If it needs an exponent above $M$, overflow occurs.

\medskip
We say the system is \newterm{parametrized}\index{parametrized} by
$(\beta, p, m, M)$.
Changing any one of the four gives a new floating-point system with different
capabilities and different limitations.

\begin{Exercise}[title={Reading the parameters}, label={ex:fp-parameter}]
  A toy floating-point system has parameters $(\beta, p, m, M) = (10, 4, -3, 4)$.
  \begin{enumerate}
    \item What base does this system use?
    \item How many significant digits does it carry?
    \item What are the smallest and largest allowed exponents?
    \item Can this system represent $12{,}345$?
          Try writing $12{,}345$ in the form $0.d_1d_2d_3d_4 \times 10^e$ and
          check whether the exponent is within bounds.
    \item Can it represent $0.00001$?
          Show your reasoning.
    \item Write two numbers that \emph{can} be represented exactly, and one
          that causes overflow.
  \end{enumerate}
\end{Exercise}
\begin{Answer}[ref={ex:fp-parameter}]
  \begin{enumerate}
    \item Base $10$.
    \item $4$ significant digits.
    \item Smallest: $m = -3$;\; largest: $M = 4$.
    \item $12{,}345 = 0.12345 \times 10^5$.
          The exponent $5$ exceeds $M = 4$, so this causes \textbf{overflow}.
          The system would need to round to $0.1235 \times 10^5 = 12{,}350$
          and even then the exponent is out of range.
    \item $0.00001 = 0.1000 \times 10^{-4}$.
          The exponent $-4$ is below $m = -3$, so this causes
          \textbf{underflow}.
    \item Any number in the range $10^{-3}$ to $(1-10^{-4}) \times 10^4$
          works, for example $0.5000 \times 10^2 = 50.00$ or
          $0.9999 \times 10^{4} = 9999$.
          Overflow example: any number with magnitude above $9999$, such
          as $100{,}000$.
  \end{enumerate}
\end{Answer}

\subsection{Connecting to Binary Floating-Point}

All of the above applies to a computer's binary floating-point system,
with $\beta = 2$ instead of $\beta = 10$.
The significand digits $d_1, d_2, \ldots, d_p$ are now bits: each one is either
$0$ or $1$.
The exponent is a power of $2$.

So a binary float looks like:
\[
  \pm \;(0.d_1\, d_2\, \cdots\, d_p)_2 \;\times\; 2^e,
  \qquad d_i \in \{0,1\}, \quad m \leq e \leq M
\]

The specific choices of $p$, $m$, and $M$ define how much memory is used and
how accurately numbers are stored.
The international standard that fixes these choices for every modern computer
is called \newterm{IEEE 754}\index{IEEE 754}, and it is the subject of the next
section.






